\section{Mikroskobik Kararlılık: Kiral Yük, Coulomb Mekanizması ve Elektronun Düşmemesi}
\label{sec:mikroskopik_kararlilik}

Klasik elektrodinamikte ivmelenen bir elektron Larmor formülü uyarınca enerji yayar:
\begin{equation}
    P = \frac{e^2 a^2}{6\pi \varepsilon_0 v_c^3}.
    \label{eq:larmor_klasik}
\end{equation}
Bu enerji kaybı nedeniyle elektronun $\sim 10^{-11}$ saniye içinde protonun üzerine sarmal çizerek çökmesi gerekirdi. Kuantum mekaniği bunu dalga fonksiyonunun belirsizlik ilkesi ($\Delta x \Delta p \ge \hbar/2$) ile postüle eder ancak fiziksel bir geometrik mekanizma sunmaz.

\subsection{Elektriksel Yükün Kiral Girdap Geometrisi}
Elektriksel yük, vakum akışkanının \textbf{Kiral Burulmasıdır (Chiral Vorticity - $\vec{\omega}$)}:
\begin{itemize}
    \item \textbf{Pozitif Yük ($+$ / Proton):} Vakum dokusunda sağ-kiral (saat yönünde) dolaşıma sahip bir girdap yutağıdır ($\vec{\omega} > 0$).
    \item \textbf{Negatif Yük ($-$ / Elektron):} Vakum dokusunda sol-kiral (saat yönünün tersine) dolaşıma sahip bir girdap yutağıdır ($\vec{\omega} < 0$).
    \item \textbf{Coulomb Çekimi (Dişli Çark Etkisi):} Zıt kiraliteli iki girdap (+ ve -) karşılaştığında, aralarındaki temas bölgesinde akış çizgileri birbirini destekleyen yönde akar. Bernoulli prensibi uyarınca ara bölgede vakum basıncı düşer ve iki parçacık birbirine doğru çekilir.
    \item \textbf{Coulomb İtmesi (Basınç Bariyeri):} Aynı kiraliteli iki girdap (+ ve + veya - ve -) karşılaştığında ise temas yüzeyinde akıntılar kafa kafaya çarpışır; yerel vakum basıncı fırlayarak parçacıkları birbirinden iter.
\end{itemize}

\subsection{Heisenberg Girdap Basıncı ve Ölçeksel Kilitlenme}
Elektronun proton üzerine çöküşü kuantum mekaniksel ve ölçek-geometrik iki temel ilkeyle imkansız kılınır:
\begin{enumerate}
    \item \textbf{Heisenberg Kuantum Kinetik Basıncı (Fiziksel Çöküş Bariyeri):} Faz uzayı diferansiyel hacmi $d\Gamma = d^3\mathbf{x} \, d^3\mathbf{p}$ topolojik bir korunum yasasıdır. Elektron çekirdeğe yaklaştıkça ($\Delta x \to 0$), Heisenberg belirsizliği gereği radyal momentum ve sıfır noktası kinetik enerjisi $E_{\mathrm{kin}} \sim \hbar^2 / (2m_e \Delta x^2)$ biçiminde $1/r^2$ ile fırlar. Buna karşılık çekici Coulomb potansiyeli yalnızca $-1/r$ ile derinleşir. Kısa mesafelerde kuantum kinetik basınç çekici potansiyeli kesin olarak yener ve sistemi $r = a_0$ Bohr yarıçapında kararlı bir taban durumuna kilitler; ışımasal çöküş dinamik olarak durdurulur.
    \item \textbf{Ölçeksel Kovaryans ve Asimptotik Çekici ($\mathbb{D}_t(r/\lambda) = 0$):} Atomun taban durumu oluştuktan sonra, evrensel ölçek büzülmesi atomun tüm yapısını ve gözlemci cetvelini eşzamanlı olarak $\lambda(t)$ ile küçültür. Ölçek-kovaryant türev altında $\mathbb{D}_t(r/\lambda) = 0$ olduğundan, bağlı sistem kendi içsel koordinatlarında mutlak bir konformal çekici (conformal attractor) olarak kararlılığını milyarlarca yıl muhafaza eder.
\end{enumerate}

\subsection{Sommerfeld İnce Yapı Sabiti ve Rölativistik Hatve Oranı}
\label{sec:alpha_main}
Kuantum elektrodinamiğinde boyutsuz elektromanyetik kuplaj sabiti olan Sommerfeld ince yapı sabiti:
\begin{equation}
    \alpha \equiv \frac{e^2}{4\pi \varepsilon_0 \hbar c} \approx \frac{1}{137.035999...},
\end{equation}
Ölçek-Zaman formalizminde iki tamamlayıcı fiziksel perspektifle analiz edilir:
\begin{enumerate}
    \item \textbf{Kinematik Spiral Hatve Oranı:} Hidrojen atomunda elektronun en iç Bohr yörünge hızı $v_{\mathrm{Bohr}} = e^2 / (4\pi \varepsilon_0 \hbar) \approx 2187\text{ km/s}$'dir. Vakum dokusundaki karakteristik içsel faz iletim tavan hızı ise $v_c = c$'dir. Bu iki hızın oranı:
          \begin{equation}
              \tan\theta_{\mathrm{hatve}} \equiv \frac{v_{\mathrm{Bohr}}}{v_c} = \frac{e^2}{4\pi \varepsilon_0 \hbar c} \equiv \alpha \approx \frac{1}{137.036},
          \end{equation}
          bağıntısını verir. Bu kinematik bağıntı, kuantum elektrodinamik kuplajın ölçek kalkülüsünde radyal faz tavanı ile orbital teğetsel hız arasındaki rölativistik hatve açısı olarak temsil edildiğini gösterir.
    \item \textbf{Kızılötesi Sınır ve Süperkritik Dirac Eşiği ($Z\alpha \to 1$):} QED'de $\alpha$ vakum polarizasyonu sebebiyle enerjiye bağlı koşan bir parametredir ($\beta(\alpha) = \alpha^2/3\pi$). Atomik bağlı durumların yaşadığı asimptotik kızılötesi (Thomson) sınırında ($Q^2 \to 0$), $\alpha \approx 1/137.036$ değeri yörünge hatve eğimini kilitler. Ağır çekirdeklerde rölativistik Dirac sınırına ($Z\alpha \to 1$) ulaşıldığında ise temel durum negatif enerjili Dirac denizine gömülerek kendiliğinden elektron-pozitron çift üretimiyle (Schwinger/Pomeranchuk mekanizması) süperkritik alanı perdeler.
\end{enumerate}

\begin{quote}
    \small
    \textbf{Kavramsal ve Fiziksel Özet (Atomik Kararlılık ve Kiral Yük):} \\
    1. Elektrik yükü, durağan vakum dokusundaki topolojik faz girdabının kiralitesini (dönel yönelimini) temsil eder; karşıt kiraliteler arasındaki faz gradyanı çekici Coulomb potansiyelini üretir. \\
    2. Elektronun çekirdeğe düşmemesi, Heisenberg belirsizlik ilkesi ($\Delta x \Delta p \ge \hbar/2$) uyarınca yarıçap küçüldükçe momentum ve sıfır noktası kinetik enerjisinin hızla artarak dışa doğru bir kuantum basınç bariyeri oluşturmasıyla güvenceye alınır. \\
    3. Büzülen gözlemci çerçevesinde yerel uzunluk cetveli de atomla eş zamanlı küçüldüğünden ($\mathbb{D}_t(r/\lambda) = 0$), bağlı sistem dinamik bir asimptotik çekici durumunda kararlı kalır.
\end{quote}

